\chapter{Objectives and Contributions}
\label{ch:objectives}

The main objective of this thesis is to investigate how structured knowledge represented by a medical knowledge graph can be combined with statistical evidence derived from patient-level data.
The problem is studied through two complementary tasks: graph reconstruction and patient-level clinical endpoint prediction.

The specific contributions of the thesis are:

\paragraph{Task A. Directed link prediction.}
A method is developed for evaluating candidate directed relations in a heterogeneous pharmacotherapy 
graph by combining graph structure with explicit empirical evidence from patient-level data. 
The graph component, implemented with a Heterogeneous Graph Transformer (HGT), provides a 
topology-constrained representation of the candidate relation and its local typed context. 
The statistical component characterizes the geometry and strength of dependence observed in 
patient-level data using coefficients derived from Hierarchical Correlation Reconstruction (HCR), 
together with additional statistical descriptors. Both representations are then combined and passed 
to a common decoder for directed link prediction. The same principle is additionally examined using 
Last-FM, a corrected music recommendation benchmark represented as a user–artist interaction graph, 
which serves as a proxy task under a substantially different graph structure and ranking objective. 
This allows us to assess whether the value of combining structural and explicit empirical 
information is specific to the pharmacotherapy setting or persists after a substantial change 
of domain.

\paragraph{Task B. Patient-level endpoint prediction.}
A patient-specific prediction framework is developed in which the shared pharmacotherapy graph provides the structural representation and patient observations provide individual input information.
Several GNN architectures and graph depths are evaluated, including the effect of residual connections on deeper message passing.
The graph representation is further complemented with an HCR-derived statistical side-channel describing dependencies between endpoints and their direct or higher-order graph neighbourhood.
Different forms of this statistical branch, including pairwise, type-specific and higher-order representations, are compared.
The resulting models are evaluated separately for clinical endpoints using metrics appropriate for imbalanced outcomes, including AUROC, Average Precision (AP) and lift.

In both cases, graph-based representation learning and explicit statistical dependence are treated as complementary sources of information rather than as competing alternatives.
